\chapter{Introduction}
\label{ch:1}
\label{sec:1}


\section{Background and Motivation}
\label{sec:1_1}


The building sector accounts for approximately 40\% of final energy consumption in the European Union and generates approximately 36\% of its energy-related greenhouse gas emissions. The 2024 recast of the Energy Performance of Buildings Directive (EPBD) further established the policy objective of transforming the existing building stock into a zero-emission building stock by 2050 \hyperlink{ref:epbd2024}{(European Parliament and Council of the European Union, 2024)}. However, the actual coverage of Energy Performance Certificate (EPC) data remains incomplete. The Netherlands introduced its EPC scheme in 2008, but by the end of 2019, only approximately 3.8 million residential buildings had a valid, registered EPC, representing around 48\% of the residential building stock \hyperlink{ref:caepbd2020}{(Concerted Action EPBD, 2020)}. Large-scale model-based estimation can therefore complement formal energy certification for building-stock analysis and preliminary retrofit prioritisation, although its use should be clearly distinguished from legally issued EPC certification.


Urban Building Energy Modelling (UBEM) provides a bottom-up framework for analysing the operational energy use of tens to thousands of buildings at neighbourhood or city scale \hyperlink{ref:reinhart2016}{(Reinhart \& Cerezo Davila, 2016)}. Because collecting detailed information on the geometry, construction, systems, and use of individual buildings is prohibitively costly, many UBEM workflows employ an archetype-based approach. Buildings with similar characteristics are grouped into categories, and each building inherits the thermal and system parameters associated with its archetype. TABULA is a representative European building-typology framework that defines residential typologies for 20 European countries according to building size, construction period, and other country-specific parameters \hyperlink{ref:loga2016}{(Loga et al., 2016)}. In the TABULA-NL configuration adopted in this study, building type and the construction period derived from construction year jointly determine the archetype cell, while floor count serves as an additional feature in the subsequent energy classification. Consequently, the large-scale application of archetype-based UBEM requires not only building geometry but also semantic attributes that are frequently missing, including building type, construction year, and floor count.


Previous research has shown that archetype lookup provides a consistent representation of the building stock but cannot capture all building-specific variation. Using approximately 2.57 million Dutch residential EPC records, \hyperlink{ref:hettinga2023}{Hettinga et al. (2023)} evaluated large-scale energy labelling and reported an accuracy of 26\% for TABULA rule-based assignment, below the majority-label baseline of 33\%. By contrast, a random forest incorporating additional building-level and neighbourhood-level open-data features achieved 71\% accuracy. Because these results were obtained using different output-label ranges and evaluation sets, their absolute values should not be interpreted as a directly comparable model benchmark. Nevertheless, they indicate that assigning energy labels solely through typology rules has limited ability to capture variation at the individual-building level.


Recent advances in remote sensing and computer vision offer another means of obtaining missing building attributes. \hyperlink{ref:wurm2021}{Wurm et al. (2021)} used aerial imagery to infer construction type and construction period and incorporated these attributes into urban heat-demand modelling. \hyperlink{ref:mayer2023uk}{Mayer et al. (2022, 2023)} used street-view and aerial imagery to predict building energy-efficiency classes directly in the United Kingdom and Denmark, respectively. OpenFACADES, meanwhile, demonstrated the feasibility of large-scale, multi-attribute building-profile enrichment using street-view imagery (SVI) and vision-language models \hyperlink{ref:liang2025}{(Liang et al., 2025)}.


However, previous studies have primarily evaluated either image-based attribute extraction or direct image-to-energy classification in isolation, leaving attribute-mediated substitution without controlled evaluation. Specifically, it remains unclear which information is retained when building type, construction year, and floor count inferred from SVI replace reference attributes within the same TABULA-grounded workflow; which of these inferred attributes provide useful information for archetype assignment and downstream EPC classification; and what performance costs and applicability boundaries arise from this substitution. This gap motivates the present study and leads to the three research objectives presented in the following section.


\section{Research Objectives}
\label{sec:1_2}


The large-scale application of archetype-based UBEM depends on the key metadata required to map individual buildings to appropriate archetypes. Existing open geospatial datasets, however, primarily provide building geometry, while semantic attributes such as building type, construction year, and floor count are frequently missing. Within the TABULA-NL framework, building type and the construction period derived from construction year determine the archetype cell, while floor count serves as an additional feature for subsequent energy classification.


The overall objective of this study is to evaluate whether building attributes inferred from street-view imagery (SVI) can replace the reference attributes provided by administrative and three-dimensional reconstruction data within a TABULA-based building-stock enrichment workflow, and to define the conditions under which such substitution is appropriate. To achieve this overall objective, the study establishes three specific research objectives:


\begin{enumerate}
\item \textbf{Evaluate SVI-based building-attribute extraction.} Assess the performance of different visual-inference configurations in predicting building type, construction year, and floor count from SVI.
\item \textbf{Determine the task-specific utility of the extracted attributes.} Determine the contribution of each SVI-derived attribute to TABULA archetype assignment and downstream EPC classification.
\item \textbf{Quantify the consequences and boundary conditions of attribute substitution.} Quantify the performance effects of replacing the full set of reference attributes with SVI-derived attributes and define the conditions under which such substitution is appropriate.
\end{enumerate}


To support these objectives, this study constructs a three-layer, value-level traceable, TABULA-grounded semi-synthetic building dataset. The dataset integrates reference observations, SVI-derived attributes, deterministically assigned TABULA archetype parameters, and EP-Online EPC reference labels, and demonstrates its analytical use through a downstream EPC-classification case study.
